\documentclass[10pt, a4paper]{article}

% --- Packages ---
\usepackage[utf8]{inputenc}
\usepackage[T1]{fontenc}
\usepackage[german]{babel}
\usepackage{geometry}
\usepackage{xcolor}
\usepackage{titlesec}
\usepackage{hyperref}
\usepackage{fontawesome5}
\usepackage{charter}

% --- Configuration ---
\geometry{left=2.5cm, top=2cm, right=2.5cm, bottom=2.5cm}
\definecolor{primarycolor}{RGB}{0, 51, 102}
\hypersetup{colorlinks=true, urlcolor=primarycolor, linkcolor=primarycolor}

% --- Document Start ---
\begin{document}

% --- Applicant Header ---
\begin{flushleft}
    \textbf{\Large \color{primarycolor} Kartavya Niraj Dikshit} \\
    \vspace{2pt}
    \small
    Halbmondstraße 7, 91054 Erlangen \\
    \faPhone* (+49) 15510940829 \quad \faEnvelope* \href{mailto:kartavya.dikshit@gmail.com}{kartavya.dikshit@gmail.com} \\
    \faLinkedin \href{http://www.linkedin.com/in/kartavya-dikshit}{linkedin.com/in/kartavya-dikshit} \quad \faGithub \href{https://github.com/KartavyaDikshit}{github.com/KartavyaDikshit}
\end{flushleft}

\vspace{20pt}

% --- Recipient Details ---
\begin{flushleft}
    \textbf{Hacon Ingenieurgesellschaft mbH} \\
    Personalabteilung \\
    Lister Meile 2 \\
    30161 Hannover \\
\end{flushleft}

\begin{flushright}
    Erlangen, 19. März 2026
\end{flushright}

\vspace{30pt}

% --- Subject ---
\noindent \textbf{\large Bewerbung als Werkstudent (d/w/m) im Bereich Data Science / Machine Learning} \\
\textbf{Referenz: Job ID 493832}

\vspace{25pt}

% --- Salutation ---
Sehr geehrte Damen und Herren,

\vspace{10pt}

% --- Introduction ---
als Masterstudent der Data Science an der Friedrich-Alexander-Universität Erlangen-Nürnberg mit einer Spezialisierung auf KI und skalierbare Systeme verfolge ich die Arbeit von Hacon im Bereich der digitalen Mobilitätslösungen mit großer Begeisterung. Die Möglichkeit, das TPS Sales Logistics Team bei der Entwicklung präziser Verspätungsprognosen zu unterstützen, ist für mich die ideale Gelegenheit, meine theoretischen Kenntnisse in einem hochrelevanten, praxisnahen Umfeld einzubringen.

% --- Experience & Technical Fit ---
In meinem aktuellen Studium und durch meine Projekterfahrung im Bereich „Real-Time Rail Delay Prediction“ habe ich mich intensiv mit der Vorhersage von Zuglaufdaten beschäftigt. Unter Einsatz von Spatio-Temporal Graph Convolutional Networks (ST-GCN) und PyTorch konnte ich eine Vorhersagegenauigkeit von 89 \% erreichen. Diese Erfahrung in der Verarbeitung komplexer Zeitreihendaten und der Optimierung von ML-Modellen für niedrige Latenzzeiten korrespondiert direkt mit den Anforderungen Ihrer ausgeschriebenen Position. Ich bin sicher im Umgang mit Python und den gängigen Bibliotheken wie TensorFlow und scikit-learn.

% --- Soft Skills & Process ---
Durch meine bisherigen Tätigkeiten als AI-ML Engineer Intern bin ich bestens mit der Arbeit in agilen Teams (Scrum) vertraut. Ich lege großen Wert auf eine strukturierte Arbeitsweise und die nachhaltige Dokumentation entwickelter Lösungen, um eine konsistente Wissensbasis im Team zu gewährleisten – ein Aspekt, den Sie in Ihrer Ausschreibung besonders betonen. Zudem bringe ich fundierte Grundkenntnisse in der Softwareentwicklung (Docker, FastAPI) mit, was mir hilft, ML-Modelle nicht nur zu trainieren, sondern auch deren Integration in bestehende Systeme zu unterstützen.

% --- Motivation & Closing ---
Hacon als Teil der Siemens Mobility steht für die Zukunft der umweltfreundlichen Fortbewegung. Es motiviert mich sehr, durch meine Arbeit einen direkten Beitrag zur Effizienz des öffentlichen Verkehrs zu leisten. Ich freue mich darauf, meine analytischen Fähigkeiten und meine Eigenständigkeit in Ihr Team in Hannover einzubringen.

Gerne überzeuge ich Sie in einem persönlichen Gespräch von meiner Eignung für diese Position. Meine Gehaltsvorstellung liegt bei 17 € pro Stunde.

\vspace{25pt}

Mit freundlichen Grüßen,

\vspace{15pt}

Kartavya Niraj Dikshit

\end{document}